\documentclass[12pt,a4paper]{article}
\usepackage{cmap} % Makes the PDF copiable. See http://tex.stackexchange.com/a/64198/25761
\usepackage[T1]{fontenc}
\usepackage[brazil]{babel}
\usepackage[utf8]{inputenc}
\usepackage{amsmath}
\usepackage{amsfonts}
\usepackage{amssymb}
\usepackage{amsthm}
\usepackage{textcomp} % \degree
\usepackage{gensymb} % \degree
\usepackage[usenames,svgnames,dvipsnames]{xcolor}
\usepackage{hyperref}
\usepackage{multicol}
\usepackage{graphicx}
\usepackage[margin=2cm]{geometry}
\usepackage{systeme}

\hypersetup{
    colorlinks = true,
    allcolors = {blue}
}

\newcommand{\fixme}{{\color{red}(...)}}
\newcommand*\ger[1]{\operatorname{ger}\left\{#1\right\}}
\newcommand\ii{\mathrm{i}}

\newcommand*\R{\mathbb{R}}

\newcommand{\IconPc}{\includegraphics[width=1em]{computer.png}}
\newcommand{\IconCalc}{\includegraphics[width=1em]{calculator.png}}
\newcommand{\IconThink}{\includegraphics[width=1em]{pencil.png}}
\newcommand{\IconCheck}{\includegraphics[width=1em]{checkmark.png}}
\newcommand{\IconConcept}{\includegraphics[width=1em]{edit.png}}

\newlength{\SmileysLength}
\setlength{\SmileysLength}{\labelwidth}\addtolength{\SmileysLength}{\labelsep}

\newcommand{\calc}{\hspace*{-\SmileysLength}\makebox[0pt][r]{\IconCalc}%
   \hspace*{\SmileysLength}}
\newcommand{\software}{\hspace*{-\SmileysLength}\makebox[0pt][r]{\IconPc}%
   \hspace*{\SmileysLength}}
\newcommand{\teoria}{\hspace*{-\SmileysLength}\makebox[0pt][r]{\IconThink}%
   \hspace*{\SmileysLength}}
\newcommand{\conceito}{\hspace*{-\SmileysLength}\makebox[0pt][r]{\IconCheck}%
   \hspace*{\SmileysLength}}
\newcommand{\concept}{\hspace*{-\SmileysLength}\makebox[0pt][r]{\IconCheck}%
   \hspace*{\SmileysLength}}

% Loop Space / CC BY-SA-3.0 / https://tex.stackexchange.com/a/2238/25761
\newenvironment{amatrix}[1]{%
  \left[\begin{array}{@{}*{#1}{c}|c@{}}
}{%
  \end{array}\right]
}

% Loop Space / CC BY-SA-3.0 / https://tex.stackexchange.com/a/3164/25761
%--------grstep
% For denoting a Gauss' reduction step.
% Use as: \grstep{\rho_1+\rho_3} or \grstep[2\rho_5 \\ 3\rho_6]{\rho_1+\rho_3}
\newcommand{\grstep}[2][\relax]{%
   \ensuremath{\mathrel{
       {\mathop{\longrightarrow}\limits^{#2\mathstrut}_{
                                     \begin{subarray}{l} #1 \end{subarray}}}}}}
\newcommand{\swap}{\leftrightarrow}

\newcommand*\tipo{3ª Lista de Exercícios}
%\newcommand*\turma{...}
\newcommand*\disciplina{ALI0001}
\newcommand*\eu{Helder G. G. de Lima}
\newcommand*\data{\today}

\author{\eu}
\title{\tipo - \disciplina}
\date{\data}

\begin{document}

\begin{center}
\includegraphics[width=9.0cm]{marca} \\
\textbf{\tipo\ (\disciplina)} \\
Prof. \eu\footnote{
Este é um material de acesso livre distribuído sob os termos da licença \href{https://creativecommons.org/licenses/by-sa/4.0/deed.pt_BR}{Creative Commons BY-SA 4.0}}
\end{center}

\section*{Legenda}
\begin{multicols}{4}
\begin{itemize}
\item[] \hspace*{\SmileysLength} \calc \hspace*{-\SmileysLength} Cálculos
\item[] \hspace*{\SmileysLength} \conceito \hspace*{-\SmileysLength} Conceitos
\item[] \hspace*{\SmileysLength} \teoria \hspace*{-\SmileysLength} Teoria
%\item[] \hspace*{\SmileysLength} \software \hspace*{-\SmileysLength} Software
\end{itemize}
\end{multicols}

\section*{Questões}

\begin{enumerate}
\item \conceito
Explique todas as afirmações verdadeiras a seguir, e dê contraexemplos para as demais.
\begin{enumerate}
\item A interseção de dois subespaços vetoriais de $V$ nunca é vazia.
\item Se $V= \ger{ u_1, u_2, u_3 }$ então $\{ u_1, u_2, u_3 \}$ é uma base de $V$
\item Se $\{u, v\}$ é linearmente dependente (L.D.) então um dos vetores é múltiplo do outro.
\item Se $u \neq 0$ então os vetores $u$ e $-u$ são linearmente independentes (L.I.).
\item Se $\{ u, v \}$ é L.I. e $\{ v, w \}$ é L.I. então $\{ u, v, w \}$ é L.I..
\item Se $\{ u \}$ é L.I. e $\{ v \}$ é L.I. então $\{ u, v \}$ é L.I..
\item É possível que $u \in \ger{v,w}$ sem que ocorra $u \in \ger{v}$ nem $u \in \ger{w}$.
\end{enumerate}

\item \calc É verdade que $\R^4 = \ger{ (1,-1,0,0), (0,0,1,-1), (0,2,1,0), (0,0,1,1)}$? Por que?

\item \calc Encontre um conjunto de vetores geradores de um dos seguintes subespaços de $M_{3 \times 3}(\R)$:
\begin{enumerate}
\item $D = \{ X \in M_{3 \times 3}(\R) \mid X \text{ é diagonal } \}$
\item $W = \{ A \in M_{3 \times 3}(\R) \mid A \text{ é antissimétrica } \}$
\item $U = \{ B \in M_{3 \times 3}(\R) \mid B \text{ é simétrica } \}$
\end{enumerate}

\item \calc Verifique se os conjuntos $\{v_1, \ldots, v_n \}$ a seguir são formados por vetores linearmente independentes (L. I.) ou linearmente dependentes (L. D.):
\begin{enumerate}
\item $v_1=(0,3,-9)$, $v_2=(1,4,-10)$, $v_3=(2,5,-11)$ em $\R^3$.
\item $v_1=(0,3,-9)$, $v_2=(1,4,-10)$ em $\R^3$.
\item $v_1 = \begin{bmatrix}
0 & 2 \\
2 & 0
\end{bmatrix}$, $v_2 = \begin{bmatrix}
1 & 0 \\
0 & 4
\end{bmatrix}$, $v_3 = \begin{bmatrix}
1 & 0 \\
0 & 0
\end{bmatrix}$, $v_4 = \begin{bmatrix}
0 & 1 \\
0 & 0
\end{bmatrix}$, $v_5 = \begin{bmatrix}
1 & 2 \\
3 & 4
\end{bmatrix}$ em $M_{2\times 2}(\R)$.
\item $v_1 = x^3 + x^2 + x + 1$, $v_2 = x^2 + x + 1$, $v_3 = x + 1$, $v_4 = 1$ no espaço de polinômios $P_4$.
\end{enumerate}

\item \calc Encontre uma base e a dimensão de cada um dos seguintes espaços vetoriais:
\begin{enumerate}
\item $V = \R^4 = \{ (x,y,z,w) \mid x \in \R, y \in \R, z \in \R, w \in \R \}$.
\item $V = M_{2 \times 3}(\R) = \left\{ \begin{bmatrix}
a & b & c \\p & q & r
\end{bmatrix} \mid a \in \R, b \in \R, c \in \R, p \in \R, q \in \R, r \in \R \right\}$.
\item $V = P_5 = \{ p(x) \in \mathcal{F}(\R) \mid p(x) \text{ é um polinômio de grau} \leq 5 \}$.
\end{enumerate}

\item \calc Encontre uma base e a dimensão do subespaço $W$ de $V$ nos seguintes casos:
\begin{enumerate}
\item $V = \R^3$ e $W = \{ (x,y,z) \mid z = 0 \}$ (o plano horizontal que passa pela origem).
\item $V = M_{2 \times 2}(\R)$ e $W = \{ S \in V \mid S^T = S \}$ (as matrizes simétricas $2 \times 2$)
\item $V = \mathcal{F}(\R)$ e $W = \{ p(x) = ax^2+bx+c \in P_2 \mid p^\prime(6) = 0 \}$ (geometricamente, estas funções correspondem a parábolas com vértice em $(6, f(6))$).
\item $V = \R^4$ e $W = \{ (x,y,z,w) \mid x + 2w = 0, -y = -2w, w=z/2 \}$.
\end{enumerate}

\item \calc Encontre um subconjunto de $S = \{ (1,1,1), (1,1,0), (1,0,1), (0,1,1) \}$ que seja base de $\R^3$.

\item \teoria Mostre que $\{ 0, v_1, \ldots, v_n \}$ é um conjunto de vetores linearmente dependente, isto é, todo conjunto de vetores contendo o vetor nulo é linearmente dependente.

\item \teoria Prove que se $W = \ger{ v_1, \ldots, v_n}$ e $z \in W$ então $W = \ger{ v_1, \ldots, v_n, z }$.

\item \teoria Prove que se $S_1 = \{ u_1, u_2, u_3 \}$ é um conjunto de vetores linearmente independentes, então $S_2 = \{ u_1, u_2 \}$ também é um conjunto de vetores linearmente independentes.

\item \calc Dada uma matriz $A$, denote por $L(A)$ o espaço linha (subespaço gerado pelas linhas de $A$), por $C(A)$ o espaço coluna (subespaço gerado pelas colunas de $A$) e por $N(A)$ o núcleo de $A$ (o espaço das matrizes coluna $X$ que satisfazem $AX = 0$). Encontre uma base e a dimensão de $L(A)$, $C(A)$ e $N(A)$ para as seguintes matrizes:
\begin{enumerate}
\begin{multicols}{2}
\item $A = \begin{bmatrix}
1 & -1 \\
1 & 1
\end{bmatrix}$

\item $A = \begin{bmatrix}
1 & 2 & 3 & 4
\end{bmatrix}$

\item $A = \begin{bmatrix}
1 & 0 & 1 & 0 \\
0 & 1 & 0 & -1
\end{bmatrix}$

\item $A = \begin{bmatrix}
1 & 0 & 2 & 0 & 3 \\
2 & 0 & 4 & 0 & 6
\end{bmatrix}^T$

\end{multicols}
\end{enumerate}

\item \calc Se $U$ e $V$ são subespaços de um espaço vetorial $W$, denota-se o conjunto de todas as somas de elementos de $U$ com elementos de $V$ por $ U + V = \{ x + y \in W \mid x \in U \text{ e } y \in V \}$.
Esta soma é chamada de \textbf{soma direta}, e denotada por $U \oplus V$, se ocorrer $U \cap V = \{ 0 \}$.

Encontre $U+V$ para os seguintes subespaços $U$ e $V$ e verifique se é uma soma direta:
\begin{enumerate}
\item Em $W = \R^4$, os subespaços $U = \{ (x,y,z,w) \in \R^4 \mid x + y = 0 = z + w \}$ e $V = \{ (a,b,c,d) \in \R^4 \mid b = 2c \}$.

\item Os subespaços $U = \{ M \in M_{3 \times 3}(\R) \mid M^T = -M \}$ das matrizes antissimétricas e $V = \{ N \in M_{3 \times 3}(\R) \mid N^T = N \}$ das matrizes simétricas de $W = M_{3 \times 3} (\R)$.

\item Os subespaços $U = \{ q(x)=b_0 + b_1x+b_2x^2+b_3x^3 \mid q(-x) = -q(x), \forall x \in \R \}$ e $V = \{ p(x)=a_0 + a_1x+a_2x^2+a_3x^3 \mid p(-x) = p(x), \forall x \in \R \}$ do espaço vetorial $P_3$.
\end{enumerate}

\item \calc Encontre uma base e a dimensão de cada uma das somas $U+V$ do exercício anterior.

\item \teoria Mostre que se $u$, $v$ e $w$ são vetores L. I. então os vetores $u+v$, $u-v$ e $u-w$ são L. I.

\item \teoria Mostre que se $B = \{u, v\}$ é uma base de $V$ então $V = \ger{u} \oplus \ger{v}$.

\item \calc Sejam $U = \ger{ (1,0,0), (2,1,1) }$ e $V = \ger{ (0,-1,0), (0,0,1) }$ subespaços de $\R^3$.
\begin{enumerate}
\item Encontre uma base e a dimensão de $U \cap V$.
\item Verifique se $\R^3 = U + V$ e se $\R^3 = U \oplus V$.
\end{enumerate}

\item \teoria Mostre que se $B = \{u, v\}$ é uma base de $V$ então $B^\prime = \{ u+v, u-v \}$ também é base.

\end{enumerate}

\newpage
\section*{Respostas}

\begin{enumerate}
\item
\begin{enumerate}
\item \textbf{Verdadeira}. Como todo subespaço vetorial deve conter o vetor nulo, a interseção de dois subespaços vetoriais quaisquer tem pelo menos um elemento, ou seja, não é vazia.
\item \textbf{Falsa}. Mesmo que $V$ seja o espaço gerado por ${ u_1, u_2, u_3 }$, pode ser que um destes vetores seja combinação linear dos demais. Neste caso, os vetores são linearmente dependentes, e não são uma base de $V$. Por exemplo: $\R^2 = \ger{(1,0),(0,1),(2,3)}$, mas $(1,1) = 2(1,0)+3(0,1)$, ou seja, estes três vetores geram $\R^2$ mas não são uma base.
\item \textbf{Verdadeira}. Se $\{u, v\}$ é linearmente dependente, então existem $a,b \in \R$ tais que $au+bv = 0$, sem que ambos os coeficientes precisem ser nulos. Se, por exemplo, $a \neq 0$, então pode-se escrever $u = (-\frac{b}{a}) v$. E se $b \neq 0$, então $v = (-\frac{a}{b}) u$. Em ambos os casos, um dos vetores pode ser visto como um múltiplo do outro.

\item \textbf{Falso}. Os vetores $u$ e $-u$ nunca são linearmente independentes, pois um deles é múltiplo do outro, e pode-se escrever $au + b(-u) = 0$ com os escalares não nulos $a=1$ e $b=1$.

\item \textbf{Falso}. Mesmo que $\{u, v\}$ sejam L.I. e que $\{ v, w \}$ seja L.I., pode acontecer de $\{u, v, w \}$ ser L.D.. Por exemplo, em $\R^2$, se $u = (1,0)$, $v = (0,1)$ e $w = (1, 1)$ então $w = u + v$, ou seja, há uma combinação linear $1u+1v+(-1)w$ que resulta no vetor nulo sem que todos os coeficientes sejam nulos.
\item \textbf{Falsa}. Apesar de, separadamente, os vetores serem L.I., pode ocorrer que o conjunto contendo ambos seja L.D. Por exemplo, em $\R^3$, o conjunto $\{ (1,0,0) \}$ é linearmente independente, e o conjunto $\{ (5,0,0) \}$ também. No entanto, como $(5,0,0) = 5 (1,0,0)$, o conjunto $\{ (1,0,0), (5,0,0) \}$ é linearmente dependente.
\item \textbf{Verdadeira}. Em $\R^2$, considerando-se $u=(2,3)$, $v = (1,0)$ e $w = (0,1)$, observa-se que $u = 2(1,0)+3(0,1) = 2u+3v \in \ger{v,w}$, mas $u \not \in \ger{v}$ e $u \not \in \ger{w}$.
\end{enumerate}

\item Para que $\R^4 = \ger{u,v,x,z}$, com $u=(1,-1,0,0)$, $v=(0,0,1,-1)$, $x=(0,2,1,0)$ e $z=(0,0,1,1)$ é necessário e suficiente que qualquer $(a,b,c,d)$ possa ser escrito como uma combinação linear dos vetores $u,v,x,z$, isto é, que dado $(a,b,c,d) \in \R^4$ exista alguma solução para a equação
\[
(a,b,c,d) = t_1 (1,-1,0,0) + t_2 (0,0,1,-1)+ t_3 (0,2,1,0)+ t_4 (0,0,1,1),
\]
que é equivalente a um sistema linear nas variáveis $t_1, t_2, t_3$ e $t_4$:

\systeme[t_1,t_2,t_3,t_4]{
 t_1                    = a,
-t_1       + 2t_3       = b,
       t_2 +  t_3 + t_4 = c,
      -t_2        + t_4 = d
}

ou ainda, à equação matricial
\[
\begin{bmatrix}
 1 &  0 & 0 & 0 \\
-1 &  0 & 2 & 0 \\
 0 &  1 & 1 & 1 \\
 0 & -1 & 0 & 1
\end{bmatrix}
\begin{bmatrix}
t_1 \\
t_2 \\
t_3 \\
t_4
\end{bmatrix}
=
\begin{bmatrix}
a \\
b \\
c \\
d
\end{bmatrix}.
\]
Note que esta matriz poderia ser construída diretamente a partir dos vetores, copiando-se as coordenadas de $u$, $v$, $x$ e $z$ para as colunas da matriz. Como a matriz é quadrada, o sistema terá solução se, e somente se, ela for inversível, o que será verdade se o determinante da matriz for diferente de zero. Calculando-o pela primeira linha, obtém-se:
\[
\begin{vmatrix}
 1 &  0 & 0 & 0 \\
-1 &  0 & 2 & 0 \\
 0 &  1 & 1 & 1 \\
 0 & -1 & 0 & 1
\end{vmatrix}
=
1 \cdot
\begin{vmatrix}
 0 & 2 & 0 \\
 1 & 1 & 1 \\
-1 & 0 & 1
\end{vmatrix}
=
-2 \cdot
\begin{vmatrix}
 1 & 1 \\
-1 & 1
\end{vmatrix}
= -2(1 \cdot 1 - (-1) \cdot 1)
=-4 \neq 0.
\]
Então, de fato, qualquer $(a,b,c,d)$ é combinação linear dos quatro vetores dados, isto é, $R^4 = \ger{u,v,x,z}$.
\item
\begin{enumerate}
\item Para toda matriz diagonal $X \in M_{3 \times 3}(\R)$ tem-se
\[
X
= \begin{bmatrix}
a & 0 & 0 \\
0 & b & 0 \\
0 & 0 & c
\end{bmatrix}
=
a
\begin{bmatrix}
1 & 0 & 0 \\
0 & 0 & 0 \\
0 & 0 & 0
\end{bmatrix}
+b
\begin{bmatrix}
0 & 0 & 0 \\
0 & 1 & 0 \\
0 & 0 & 0
\end{bmatrix}
+c
\begin{bmatrix}
0 & 0 & 0 \\
0 & 0 & 0 \\
0 & 0 & c
\end{bmatrix},
\]
com $a,b,c \in \R$. Então estas três matrizes geram $D$.

\item Para toda matriz antissimétrica $A \in M_{3 \times 3}(\R)$ tem-se
\[
A
= \begin{bmatrix}
 0 & a & b \\
-a & 0 & c \\
-b & c & 0
\end{bmatrix}
=
a
\begin{bmatrix}
 0 & 1 & 0 \\
-1 & 0 & 0 \\
 0 & 0 & 0
\end{bmatrix}
+b
\begin{bmatrix}
 0 & 0 & 1 \\
 0 & 0 & 0 \\
-1 & 0 & 0
\end{bmatrix}
+c
\begin{bmatrix}
0 &  0 & 0 \\
0 &  0 & 1 \\
0 & -1 & 0
\end{bmatrix},
\]
com $a,b,c \in \R$. Então estas três matrizes geram $W$.

\item Para toda matriz simétrica $B \in M_{3 \times 3}(\R)$ tem-se
\begin{align*}
B
= \begin{bmatrix}
a & b & c \\
b & d & e \\
c & e & f
\end{bmatrix}
&=
a
\begin{bmatrix}
1 & 0 & 0 \\
0 & 0 & 0 \\
0 & 0 & 0
\end{bmatrix}
+b
\begin{bmatrix}
0 & 1 & 0 \\
1 & 0 & 0 \\
0 & 0 & 0
\end{bmatrix}
+c
\begin{bmatrix}
0 & 0 & 1 \\
0 & 0 & 0 \\
1 & 0 & 0
\end{bmatrix}
+d
\begin{bmatrix}
0 & 0 & 0 \\
0 & 1 & 0 \\
0 & 0 & 0
\end{bmatrix}\\
&+e
\begin{bmatrix}
0 & 0 & 0 \\
0 & 0 & 1 \\
0 & 1 & 0
\end{bmatrix}
+f
\begin{bmatrix}
0 & 0 & 0 \\
0 & 0 & 0 \\
0 & 0 & 1
\end{bmatrix},
\end{align*}
com $a,b,c,d,e,f \in \R$. Então estas 6 matrizes geram $U$.
\end{enumerate}

\item
\begin{enumerate}
\item $v_1=(0,3,-9)$, $v_2=(1,4,-10)$, $v_3=(2,5,-11)$ em $\R^3$.

Suponha que $c_1$, $c_2$ e $c_3$ sejam escalares tais que $c_1v_1 + c_2v_2 + c_3v_3 = 0$, isto é, que
\[
c_1(0,3,-9) + c_2(1,4,-10) + c_3 (2,5,-11)
= (0, 0, 0)
\]
ou ainda,
\[(c_2 + 2c_3, 3c_1 + 4c_2 + 5c_3, -9c_1 - 10c_2 - 11c_3)
= (0, 0, 0)
\]
Então $c_1, c_2, c_3$ são soluções do sistema linear homogêneo
\systeme{
c_2 + 2c_3 =0,
3c_1 + 4c_2 + 5c_3=0,
-9c_1 - 10c_2 - 11c_3=0
}
Como o número de equações é o mesmo do número de incógnitas, pode-se determinar se $c_1=c_2=c_3=0$ é ou não a única solução calculando o determinante da matriz associada ao sistema:
\[
\begin{vmatrix}
 0 &   1 &   2 \\
 3 &   4 &   5 \\
-9 & -10 & -11
\end{vmatrix}
=
(-1)
\begin{vmatrix}
 3 &   5 \\
-9 & -11
\end{vmatrix}
+2
\begin{vmatrix}
 3 &   4 \\
-9 & -10
\end{vmatrix}
=
(-1) \cdot 12
+2 \cdot 6
= 0
\]
Como o determinante é zero, existem infinitas soluções não triviais e os vetores dados  são linearmente dependentes.




\item $v_1=(0,3,-9)$, $v_2=(1,4,-10)$ em $\R^3$.
Pelo exercício 1, se $v_1$ e $v_2$ fossem linearmente dependentes, um deles seria múltiplo do outro. Mas estes dois vetores não são múltiplos, logo, só podem ser linearmente independentes.


\item A condição $c_1 v_1 + c_2 v_2 + c_3 v_3 + c_4 v_4 + c_5 v_5 = 0$ equivale às seguintes condições:
\begin{align*}
& c_1
\begin{bmatrix}
0 & 2 \\
2 & 0
\end{bmatrix}
c_2
\begin{bmatrix}
1 & 0 \\
0 & 4
\end{bmatrix}
c_3
\begin{bmatrix}
1 & 0 \\
0 & 0
\end{bmatrix}
c_4
\begin{bmatrix}
0 & 1 \\
0 & 0
\end{bmatrix}
c_5
\begin{bmatrix}
1 & 2 \\
3 & 4
\end{bmatrix}
=\begin{bmatrix}
0 & 0 \\
0 & 0
\end{bmatrix} \\
&
\Leftrightarrow
\systeme{
c_2+c_3+c_5=0,
2c_1+c_4+2c_5=0,
2c_1+3c_5=0,
4c_2+4c_5=0
}
\Leftrightarrow
\begin{bmatrix}
0 & 1 & 1 & 0 & 1 \\
2 & 0 & 0 & 1 & 2 \\
2 & 0 & 0 & 0 & 3 \\
0 & 4 & 0 & 0 & 4
\end{bmatrix}
\cdot
\begin{bmatrix}
c_1\\
c_2\\
c_3\\
c_4\\
c_5
\end{bmatrix}
=
\begin{bmatrix}
0\\
0\\
0\\
0\\
0
\end{bmatrix}
\end{align*}
Escalonando a matriz deste sistema (ou encontrando as suas soluções por qualquer outro método) obtém-se:
\begin{align*}
& \begin{bmatrix}
0 & 1 & 1 & 0 & 1 \\
2 & 0 & 0 & 1 & 2 \\
2 & 0 & 0 & 0 & 3 \\
0 & 4 & 0 & 0 & 4
\end{bmatrix}
\grstep{ L_1 \swap L_2 }
\begin{bmatrix}
2 & 0 & 0 & 1 & 2 \\
0 & 1 & 1 & 0 & 1 \\
2 & 0 & 0 & 0 & 3 \\
0 & 4 & 0 & 0 & 4
\end{bmatrix}
\grstep[ \frac{1}{4}L_4 ]{ \frac{1}{2}L_1 }
\begin{bmatrix}
1 & 0 & 0 & 1/2 & 1 \\
0 & 1 & 1 & 0 & 1 \\
2 & 0 & 0 & 0 & 3 \\
0 & 1 & 0 & 0 & 1
\end{bmatrix}\\
\grstep[ L_4 - L_2 ]{ L_3 - 2L_1 }
& \begin{bmatrix}
1 & 0 & 0 & 1/2 & 1 \\
0 & 1 & 1 &  0 & 1 \\
0 & 0 & 0 & -1 & 1 \\
0 & 0 & -1 & 0 & 0
\end{bmatrix}
\grstep{ L_3 \swap L_4 }
\begin{bmatrix}
1 & 0 & 0 & 1/2 & 1 \\
0 & 1 & 1 &  0 & 1 \\
0 & 0 & -1 & 0 & 0 \\
0 & 0 & 0 & -1 & 1
\end{bmatrix}
\grstep[ -L_4 ]{ -L_3 }
\begin{bmatrix}
1 & 0 & 0 & 1/2 & 1 \\
0 & 1 & 1 &  0 & 1 \\
0 & 0 & 1 & 0 & 0 \\
0 & 0 & 0 & 1 & -1
\end{bmatrix}\\
\grstep[ L_1-\frac{1}{2}L_4 ]{ L_2-L_3 }
&\begin{bmatrix}
1 & 0 & 0 & 0 & 3/2 \\
0 & 1 & 0 & 0 & 1 \\
0 & 0 & 1 & 0 & 0 \\
0 & 0 & 0 & 1 & -1
\end{bmatrix}
\Rightarrow
\begin{cases}
c_1 =-c_5/2\\
c_2 = -c_5\\
c_3 = 0\\
c_4 = c_5
\end{cases}
\end{align*}
Como o sistema tem $c_5$ como uma variável livre, é possível obter o vetor nulo como combinação linear das matrizes $v_1,\ldots,v_5$ sem que todos os coeficientes sejam nulos (exemplo: $-1v_1 - 2v_2 + 0v_3 + 5v_4 + 5v_5 = 0$). Isto significa que estes vetores são linearmente dependentes.

\item A condição $a v_1 + b v_2 + c v_3 + d v_4 = 0$ equivale às seguintes condições:
\begin{align*}
& a ( x^3 + x^2 + x + 1 ) + b ( x^2 + x + 1 ) + c (x + 1) + d
= 0 x^3 + 0 x^2 + 0 x + 0 \\
& \Leftrightarrow
a x^3 + (a+b) x^2 + (a+b+c)x + (a+b+c+d) = 0 x^3 + 0 x^2 + 0 x + 0 \\
\Leftrightarrow
&
\systeme{
a=0,
a+b=0,
a+b+c=0,
a+b+c+d=0
}
\Leftrightarrow
\begin{cases}
a=0\\
b=0\\
c=0\\
d=0
\end{cases}
\end{align*}
\end{enumerate}
Como a única forma de obter o vetor nulo como combinação linear dos polinômios dados é usar todos os coeficientes iguais a zero, resulta que eles são linearmente independentes.


\item
\begin{enumerate}
\item Pode-se verificar que $B = \{ (1,0,0,0), (0,1,0,0), (0,0,1,0), (0,0,0,1) \}$ é uma base de $\R^4$ e disto resulta que $\dim( \R^4 ) = 4$.
\item $B = \left\{
\begin{bmatrix}
1 & 0 & 0 \\0 & 0 & 0
\end{bmatrix},
\begin{bmatrix}
0 & 1 & 0 \\0 & 0 & 0
\end{bmatrix},
\begin{bmatrix}
0 & 0 & 1 \\0 & 0 & 0
\end{bmatrix},
\begin{bmatrix}
0 & 0 & 0 \\1 & 0 & 0
\end{bmatrix},
\begin{bmatrix}
0 & 0 & 0 \\0 & 1 & 0
\end{bmatrix},
\begin{bmatrix}
0 & 0 & 0 \\0 & 0 & 1
\end{bmatrix}
\right\}$ é uma base de $M_{2 \times 3}(\R)$ (por que?) e consequentemente $\dim( M_{2 \times 3}(\R) ) = 6$.

\item $B = \left\{1,x,x^2,x^3,x^4,x^5 \right\}$ é uma base de $P_5$ (verifique!) e portanto $\dim( P_5 ) = 6$.
\end{enumerate}

\item
\begin{enumerate}
\item Tem-se $u \in W$ se e somente se $u = (x,y,0)$, para algm $x,y, \in \R$. Neste caso, $u = x(1,0,0) + y(0,1,0)$. Portanto os vetores $v_1 = (1,0,0)$ e $v_2 = (0,1,0)$ geram $W$. Mas estes vetores são linearmente independentes, pois se $k_1 v_1 + k_2 v_2 = 0$ então
\[
k_1(1,0,0) + k_2(1,0,0) = (k_1,k_1,0) = (0,0,0)
\]
o que significa que $k_1 = k_2 = 0$. Em outras palavras, o vetor nulo só pode ser obtido como combinação linear de $v_1$ e $v_2$ se todos os coeficientes forem nulos. Portanto, $B = \{(1,0,0),(0,1,0)\}$ é uma base de $W$, e conclui-se que $\dim(W)=2$.

\item Um vetor $u \in V$ pertence a $W$ se, e somente se
\[
u
=
\begin{bmatrix}
a & b \\ b & c
\end{bmatrix}
=
a\begin{bmatrix}
1 & 0 \\ 0 & 0
\end{bmatrix}
+b\begin{bmatrix}
0 & 1 \\ 1 & 0
\end{bmatrix}
+c\begin{bmatrix}
0 & 0 \\ 0 & 1
\end{bmatrix},
= a v_1 + bv_2 + cv_3
\]
em que $v_1
=
\begin{bmatrix}
1 & 0 \\ 0 & 0
\end{bmatrix}$,
$v_2=
\begin{bmatrix}
0 & 1 \\ 1 & 0
\end{bmatrix}$
e
$v_3=
\begin{bmatrix}
0 & 0 \\ 0 & 1
\end{bmatrix}$ e $a,b,c \in \R$. Isto significa que $W = \ger{v_1,v_2,v_3}$. Além disso, se $m_1 v_1 + m_2 v_2 + m_3v_3 = 0$, tem-se
\[
m_1
\begin{bmatrix}
1 & 0 \\ 0 & 0
\end{bmatrix}
+m_2
\begin{bmatrix}
0 & 1 \\ 1 & 0
\end{bmatrix}
+m_3
\begin{bmatrix}
0 & 0 \\ 0 & 1
\end{bmatrix}
=
\begin{bmatrix}
m_1 & m_2 \\ m_2 & m_3
\end{bmatrix}
=\begin{bmatrix}
0 & 0 \\ 0 & 0
\end{bmatrix},
\]
isto é, cada $m_i = 0$, o que significa que os vetores $v_1, v_2, v_3$ são  L.I. e $B = \{v_1, v_2, v_3\}$ é uma base de $W$. Consequentemente, $\dim(W) = 3$.

\item $V = \mathcal{F}(\R)$ e $W = \{ p(x) = ax^2+bx+c \in P_2 \mid p^\prime(6) = 0 \}$ (geometricamente, estas funções correspondem a parábolas com vértice em $(6, f(6))$).

Dado $p(x) = ax^2+bx+c \in P_2$ tem-se $p^\prime(x) = 2ax + b$. Então $p(x) \in W$ se, e somente se $2a(6) + b = 0$, ou seja, $b = -12a$. Assim, todo elemento de $W$ é da forma
\[
p(x)
= ax^2 + (-12a)x + c
= a \cdot (x^2 -12x) + c \cdot 1,
\]
com $a,c \in \R$. Em outras palavras, os polinômios $q_1 = x^2 -12x$ e $q_2 = 1$ geram $W$. Estes vetores também são linearmente independentes pois

$
c_1(x^2 -12x) + c_2(1) = 0
\Leftrightarrow
c_1x^2 -12 c_1 x + c_2 = 0
\Leftrightarrow
\begin{cases}
    c_1 = 0\\
-12 c_1 = 0\\
    c_2 = 0
\end{cases}
\Leftrightarrow
\begin{cases}
c_1 = 0\\
c_2 = 0
\end{cases}
$
Portanto, $B = \ger{1, x^2 - 12x}$ é uma base de $W$ e $\dim(W) = 2$.

\item Como
\[
\begin{cases}
x +  2w = 0\\
  -y    = -2w\\
      w = z/2
\end{cases}
\Leftrightarrow
\begin{cases}
x = -2w\\
y = 2w\\
z = 2w
\end{cases}
\]
resulta que $u \in W$ se, e somente se, $u = (-2w,2w,2w,w) = w (-2,2,2,1)$, com $w \in \R$. Em outras palavras, $W = \ger{ (-2,2,2,1) }$. Como este vetor é diferente de zero, $B = \{ (-2,2,2,1) \}$ é L.I. pois
\[
k (-2,2,2,1) = (0,0,0,0)
\Rightarrow
\begin{cases}
-2k = 0\\
2k = 0\\
2k = 0\\
k = 0
\end{cases}
\Rightarrow k = 0.
\]
Portanto $B$ é uma base para $W$ e $\dim(W) = 1$.
\end{enumerate}

\item \calc Encontre um subconjunto de $S = \{ (1,1,1), (1,1,0), (1,0,1), (0,1,1) \}$ que seja base de $\R^3$.

Como $(1,1,1) = \frac{1}{2} \left( (1,1,0) + (1,0,1) + (0,1,1) \right)$, pode-se dizer que $S$ é L.D. (pois um dos vetores é combinação linear dos demais, ou porque o sistema linear associado tem soluções não nulas) e portanto não forma uma base de $\R^3$. No entanto, se forem considerados apenas os vetores $v_1 = (1,1,0)$, $v_2 = (1,0,1)$ e $v_3 = (0,1,1)$, então o subconjunto $B = \{ v_1, v_2, v_3 \}$ será uma base de $\R^3$ pois $R^3 = \ger{ v_1, v_2, v_3 }$ e estes três vetores são L.I. De fato, qualquer $(a,b,c) \in \R^3$ pode ser escrito como combinação linear de $v_1, v_2$ e $v_3$ pois
\begin{align*}
& (a,b,c) = x(1,1,0) + y(1,0,1) + z(0,1,1)
\Leftrightarrow
\begin{cases}
x+y=a\\
x+z=b\\
y+z=c\\
\end{cases}
\Leftrightarrow
\begin{cases}
x=\frac{1}{2}( a+b-c)\\
y=\frac{1}{2}( a-b+c)\\
z=\frac{1}{2}(-a+b+c)
\end{cases}
%\\
%& \Leftrightarrow
%(a,b,c) =
%  \left(\frac{1}{2}( a+b-c)\right)(1,1,0)
%+ \left(\frac{1}{2}( a-b+c)\right)(1,0,1)
%+ \left(\frac{1}{2}(-a+b+c)\right)(0,1,1)
\end{align*}
Além disso, com base no cálculo acima, o único modo de obter $(a,b,c) = (0,0,0)$ como combinação linear de $v_1$, $v_2$ e $v_3$ é considerando todos os coeficientes $x=y=z=0$.

\item Para que um conjunto de vetores seja L.D. deve ser possível expressar o vetor nulo como combinação linear deles sem que todos os coeficientes desta combinação linear sejam zero. No caso dos vetores $0, v_1, \ldots, v_n$, pode-se usar, por exemplo, os escalares $c_0 =7, c_1 = 0, c_2 = 0, \ldots, c_n = 0$ pois $7 \cdot 0 + 0 \cdot v_1 + \ldots + 0 \cdot v_n = 0$. Como pelo menos o escalar $c_0 \neq 0$, conclui-se que os vetores são linearmente dependentes.

\item Sejam $W = \ger{ v_1, \ldots, v_n }$ e $z \in W$, isto é, $z = k_1 v_1 + \ldots + k_n v_n$, com $k_1, \ldots, k_n \in \R$. Sendo $U = \ger{ v_1, \ldots, v_n, z }$, pode-se dizer que os conjuntos $U$ e $W$ são iguais se for explicado por que todo elemento de $W$ pertence a $U$ e vice-versa ($W \subseteq U$ e $U \subseteq W$).
\begin{itemize}
\item
Se $u \in W$, então $u$ é uma combinação linear dos vetores $v_i$, isto é, $u = c_1 v_1 + \ldots + c_n v_n$. Mas então é claro que $u$ também é combinação linear dos vetores $v_i$ e de $z$ pois basta considerar um coeficiente zero para $z$, como em $u = c_1 v_1 + \ldots + c_n v_n + 0 \cdot z$. Isto quer dizer que $z \in \ger{ v_1, \ldots, v_n, z } = U$, ou seja, que $W \subseteq U$.
\item Por outro lado, todo $u \in W$ é combinação linear de $v_1, \ldots, v_n, z$, isto é,
\begin{align*}
u
& = d_1 v_1 + \ldots + d_n v_n + d_{n+1} z \\
& = d_1 v_1 + \ldots + d_n v_n + d_{n+1}(k_1 v_1 + \ldots + k_n v_n) \\
& = (d_1 + d_{n+1}k_1) v_1 + \ldots + ( d_n + d_{n+1}k_n) v_n.
\end{align*}
Portanto $u$ é uma combinação linear dos vetores $v_1, \ldots, v_n$, o que quer dizer que $u \in \ger{ v_1, \ldots, v_n } = W$, ou seja, que $U \subseteq W$.
\end{itemize}
Em resumo, ao adicionar a uma lista de vetores $v_i$ um vetor $z$ que já esteja no espaço gerado por eles, o espaço gerado por todos estes vetores continuará sendo o mesmo.

\item Suponha que $a u_1 + bu_2 = 0$. Então $a u_1 + b u_2 + 0 u_3= 0$, ou seja, tem-se uma combinação de $u_1, u_2, u_3$ que resulta no vetor nulo. Mas por hipótese, $u_1, u_2, u_3$ são L.I., então esta combinação linear precisa ter todos os coeficientes iguais a zero, isto é, $a=b=0$. Assim, mostrou-se que não há outra forma de obter o vetor nulo como combinação linear de $u_1$ e $u_2$, isto é, $u_1$ e $u_2$ são L.I.

\item
\begin{enumerate}
\item
\begin{enumerate}
\item Como o espaço linha da matriz $A$ e o espaço linha de sua forma escalonada reduzida são os mesmos, pode-se utilizar as linhas não nulas da forma reduzida como vetores geradores de $L(A)$. Para a matriz dada, o escalonamento consistirá destes passos:
\begin{align*}
\begin{bmatrix}
1 & -1 \\
1 & 1
\end{bmatrix}
\grstep{ L_2 - L_1 }
\begin{bmatrix}
1 & -1 \\
0 & 2
\end{bmatrix}
\grstep{ \frac{1}{2}L_2 }
\begin{bmatrix}
1 & -1 \\
0 & 1
\end{bmatrix}
\grstep{ L_1 + L_2 }
\begin{bmatrix}
1 & 0 \\
0 & 1
\end{bmatrix}
\end{align*}
Então $L(A) = \ger{
\begin{bmatrix}
1 & 0
\end{bmatrix},
\begin{bmatrix}
0 & 1
\end{bmatrix}
} = \{
\begin{bmatrix}
a & b
\end{bmatrix} \mid a,b, \in \R
\} = M_{1 \times 2} (\R)$. Além disso, é claro que estes dois vetores linha são L.I. pois $a
\begin{bmatrix}
1 & 0
\end{bmatrix}
+b
\begin{bmatrix}
0 & 1
\end{bmatrix}
=
\begin{bmatrix}
a & b
\end{bmatrix}$ só pode ser igual a $
\begin{bmatrix}
0 & 0
\end{bmatrix}$ se $a$ e $b$ forem ambos nulos. Em outras palavras, estes dois vetores formam uma base de $L(A)$, que portanto é um espaço vetorial de dimensão dois.

\item
Para o espaço coluna de $A$, pode-se usar como geradores as colunas de $A$ que correspondem às posições dos pivôs da forma reduzida. Especificamente, tem-se $C(A)
= \ger{\begin{bmatrix}
1\\
1
\end{bmatrix},
\begin{bmatrix}
-1\\
1
\end{bmatrix}}$. Estes dois vetores coluna são $L.I.$ pois nenhum deles é múltiplo do outro. Então eles formam uma base de $C(A)$ e $\dim(C(A)) = 2$.
\item O núcleo $N(A)$ consiste dos vetores $\begin{bmatrix}
x\\y
\end{bmatrix}$ tais que $\begin{bmatrix}
1 & -1\\
1 & 1
\end{bmatrix}\begin{bmatrix}
x\\
y
\end{bmatrix}
=
\begin{bmatrix}
0\\
0
\end{bmatrix}$, isto é, as coordenadas $x,y$ são soluções do sistema linear
\systeme{
x-y=0,
x+y=0.
}

Isto implica que $x=y=0$ e que $N(A) = \left\{ \begin{bmatrix}
0\\
0
\end{bmatrix} \right\}$. Portanto, o conjunto vazio $\emptyset$ é uma base de $N(A)$. Perceba que não faria sentido que a base tivesse algum vetor não nulo, pois senão ela geraria um espaço maior do que o espaço nulo; e também não faz sentido considerar apenas o vetor nulo como uma base, já que ele não seria L.I.. Logo a dimensão de $N(A)$ é zero (não há vetores na base).
\end{enumerate}

\item Note que $A = \begin{bmatrix}
1 & 2 & 3 & 4
\end{bmatrix}$ está na forma escalonada reduzida por linhas, então:
\begin{enumerate}
\item $L(A) = \ger{ \begin{bmatrix}
1 & 2 & 3 & 4
\end{bmatrix} }$ e como este conjunto é L.I., $\dim(L(A)) = 1$

\item
$C(A) = \ger{ \begin{bmatrix}
1
\end{bmatrix} }$ e como o conjunto é L.I. $\dim(C(A)) = 1$

\item $N(A) = \left\{ \begin{bmatrix}
-2y-3z-4w\\y\\z\\w
\end{bmatrix} \mid y,z,w\in \R
\right\}
=
\ger{ \begin{bmatrix}
-2\\1\\0\\0
\end{bmatrix},
\begin{bmatrix}
-3\\0\\1\\0
\end{bmatrix},
\begin{bmatrix}
-4\\0\\0\\1
\end{bmatrix}  }$ e como os três vetores são L.I., $\dim(L(A)) = 3$

\end{enumerate}

\item A matriz $A = \begin{bmatrix}
1 & 0 & 1 & 0 \\
0 & 1 & 0 & -1
\end{bmatrix}$ está na forma escalonada reduzida então:
\begin{enumerate}
\item $L(A) = \ger{ \begin{bmatrix}
1 & 0 & 1 & 0
\end{bmatrix}, \begin{bmatrix}
0 & 1 & 0 & -1
\end{bmatrix} }$ e $\dim(L(A)) = 2$ (os vetores são L.I.)

\item $C(A) = \ger{ \begin{bmatrix}
1 \\ 0
\end{bmatrix}, \begin{bmatrix}
0 \\ 1
\end{bmatrix} }$ e $\dim(C(A)) = 2$.

\item $N(A) = \left\{ \begin{bmatrix}
-z\\w\\z\\w
\end{bmatrix} \mid y,z,w\in \R
\right\}
=
\ger{ \begin{bmatrix}
-1\\0\\1\\0
\end{bmatrix},
\begin{bmatrix}
0\\1\\0\\1
\end{bmatrix} }$ e $\dim(L(A)) = 2$.
\end{enumerate}


\item A forma escalonada reduzida de $A = \begin{bmatrix}
1 & 2\\
0 & 0\\
2 & 4\\
0 & 0\\
3 & 6
\end{bmatrix}$ é $\begin{bmatrix}
1 & 2\\
0 & 0\\
0 & 0\\
0 & 0\\
0 & 0
\end{bmatrix}$, então
\begin{enumerate}
\item $L(A) = \ger{ \begin{bmatrix}
1 & 2
\end{bmatrix} }$ e $\dim(L(A)) = 1$.

\item $C(A) = \ger{ \begin{bmatrix}
1\\
0\\
0\\
0\\
0
\end{bmatrix} }$ e $\dim(C(A)) = 1$.

\item $N(A) = \left\{ \begin{bmatrix}
-2y\\y
\end{bmatrix} \mid y\in \R
\right\}
=
\ger{ \begin{bmatrix}
-2\\1
\end{bmatrix} }$ e $\dim(L(A)) = 1$.
\end{enumerate}

\end{enumerate}
\item \fixme
\item \fixme
\item Suponha que $a(u+v) + b(u-v) + c(u-w) = 0$. Então $(a+b+c)u+(a-b)v + (-c)w = 0$ e como $u,v,w$ são L.I. esta combinação linear só resulta no vetor nulo se os três coeficientes forem zero, isto é, se
\[
\begin{cases}
a+b+c=0\\
a-b=0\\
-c=0\\
\end{cases}
\]
Mas a única solução deste sistema é $a=b=c=0$, o que implica que vetores $u+v$, $u-v$ e $u-w$ são L.I..


\item \fixme
\item \fixme
\item \fixme
\end{enumerate}
\end{document}
